Un rotor de inercia \(J=2\,\mathrm{kg\,m^2}\) arranca desde reposo con un par motor constante de \(1\,\mathrm{N\,m}\). En un ensayo sintético alcanza un régimen estable de \(4\,\mathrm{rad/s}\). El par resistente se opone al giro, pero no se conoce su dependencia de la velocidad. El cambio de momento angular por unidad de tiempo iguala el par neto.

Formule un modelo de arranque escogiendo y justificando una ley resistente sencilla. Ajuste lo que permita el dato y prediga cuándo se alcanza la mitad de la velocidad final. Determine si otra ley razonable compatible con el mismo ensayo puede cambiar esa predicción. Proponga una medición adicional que ayude a decidir entre ellas.
